\section{Open questions}\label{sec:program}

The difference-side approach yields both constructions and complete
exclusions on specified domains. It produces two infinite families of
$(8,7)$ arrays, classifies the $(9,7)$ arrays supported on primes at most
$19$, and closes particular full-magic fibres. The remaining questions
are arithmetic questions about the interaction, rather than about the
seven-line form itself.

First, the four-prime reductions of Section~\ref{sec:frontier} leave fixed
core fibres and unbalanced branches beyond the $S_{19}$ domain. Deciding
these branches uniformly as the two additional primes vary would extend
the finite-support results substantially. Second, fixed-squareclass
finiteness is non-effective: an effective bound or a terminating
enumeration for an arbitrary prescribed squareclass remains to be found.
Third, the model of Section~\ref{sec:surface} separates the generic cone
map from the rational-square condition needed at fixed difference. Its
other maps, covering curves and arithmetic fibres remain possible tools;
the $S_5$ result concerns the displayed map only. Likewise, the finite-class
countercertificate leaves methods using representatives, heights and
valuations available.

The companion repository records three further continuations: a residual
cover over $\Q(\sqrt{218})$, an odd lift tower, and a marked class in a
degree-$90$ $S$-class group. Their exact unresolved questions and supporting
records are separated from the principal theorems in
Appendix~\ref{app:supplementary}. None decides the general magic-square
problem. A complete two-term $S$-unit list gives a finite procedure for
one fixed support, as Section~\ref{sec:smoothsupport} demonstrates; it does
not by itself decide the problem over all supports.
